\begin{landscape}
\pdfbookmark[2]{Erste Rechnung, [4]}{pallas-III-fragment4}
\begin{center}[4.]\end{center}
\begingroup
\makeatletter\tagsleft@true\let\veqno\@@leqno\makeatother
\setlength{\lineskip}{4pt}\setlength{\lineskiplimit}{4pt}
\setlength{\arraycolsep}{3pt}
\newcommand{\IIIfourDayComma}{\mathord{\rlap{\(^{\mathrm d}\)}{,}\kern.2em}}
\newcommand{\IIIfourVStroke}{\rlap{\kern.18em\raisebox{.6ex}{\rotatebox{25}{\rule{.25em}{.2pt}}}}}
\newcommand{\IIIfourPlusMark}{\mathord{+}\text{\rlap{\kern-.1em\raisebox{1.4ex}{\tiny .}}}}
\newcommand{\IIIfourLowComma}{\mathord{\text{\rlap{\kern.1em\raisebox{-.7ex}{.}},}}}
\newcommand{\IIIfourLowC}{\mathord{\underset{\smash{\raisebox{.5ex}{\tiny .}}}{c}}}
\noindent\hspace*{\parindent}[Es handelt sich nun darum, die Werthe der Normalelemente zu bestimmen, welche in Verbindung mit diesen Störungen die Oppositionen am besten darstellen, und es soll zu diesem Zwecke im Grossen und Ganzen der in den Art. 5--6 der Disquisitio de elementis ellipticis Palladis (vgl. auch dort das Beispiel in Art. 8) auseinander gesetzten Methode benutzt werden, wo (Band VI, Seite 4) die Beobachtungen der sechs Oppositionen, wie folgt, angegeben sind:
\begin{center}
\begingroup\setlength{\tabcolsep}{8pt}
\begin{tabular}{@{}c|r|r|r@{}}
Opposition & \(1803+\) & heliocentrische Länge \(=\alpha\) & geocentrische Breite \(=\beta\)\\
\hline
I & \(181\IIIfourDayComma019120\) & \(277^0\,39'\,24'',0\) & \(+46^0\,26'\,36'',0\)\\
II & \(608,207257\) & \(337\quad0\quad36,1\) & \(+15\quad1\quad49,8\)\\
III & \(1064,468796\) & \(67\quad20\quad42,9\) & \(-54\quad30\quad54,9\)\\
IV & \(1585,609502\) & \(223\quad37\quad27,7\) & \(\IIIfourPlusMark42\quad11\quad25,6\)\\
V\IIIfourVStroke & \(2034,887176\) & \(304\quad2\quad59,7\) & \(+37\quad43\quad53,7\)\\
VI & \(2457,673843\) & \(359\quad40\quad4,4\) & \(-7\quad22\quad10,1\)
\end{tabular}\endgroup
\end{center}
\noindent und wobei folgende Bezeichnungen eingeführt wurden:
\begin{center}
\SourceLeaf{486}
\begin{tabular}{@{}r@{\ }p{.90\linewidth}@{}}
\(\alpha={}\)&beobachtete heliocentrische Länge\\[3pt]
\(\beta={}\)&beobachtete geocentrische Breite\\[3pt]
\(\delta={}\)&heliocentrische Breite, abgeleitet aus \(\alpha\) nach der Formel \(\operatorname{tang}\delta=\operatorname{tang}i\sin(\alpha-\AscendingNode)\)\\[3pt]
\(\gamma={}\)&heliocentrische Breite, abgeleitet aus \(\beta\) nach der Formel \(\sin(\beta-\gamma)=\dfrac{R\sin\beta}{r}\).
\end{tabular}
\end{center}
\par
Als genäherte Werthe für die vier Elemente \(i\), \(\AscendingNode\), \(\varphi\), \(\tilde\omega\) bezogen auf das Äquinoctium 1803,0 wurde vorausgesetzt:
\[\begin{aligned}
i&=34^0\,37'\,47''&\qquad\AscendingNode&=172^0\,28'\,32''\\
\varphi&=14\quad12\quad51&\tilde\omega&=121\quad7\quad54.
\end{aligned}\]
\par
Aus den sechs beobachteten heliocentrischen Längen, reducirt auf das eben genannte Äquinoctium:
\[\begin{array}{@{}r@{\;}r|r|r|r|r|r@{}}
\alpha={}&277^{0}\,38'59'',17 & 336^{0}\,59'12'',67 & 67^{0}\,18'16'',89 & 223^{0}\,33'50'',20 & 303^{0}\,58'20'',58 & 359^{0}\,34'27'',29
\end{array}\]
und unter Annahme der für die sechs Zeiten mit Berücksichtigung der Störungen folgenden Werte von \(i\) und \(\AscendingNode\), nemlich:
\[\begin{array}{@{}r@{\;}r|r|r|r|r|r@{}}
\AscendingNode={}&172^{0}\,28'32'',43 & 28'49'',88 & 28'32'',77 & 28'19'',28 & 27'48'',22 & 27'45'',90
\\
i={}&34\ 37\ 46,96 & 37\ 35,85 & 37\ 12,42 & 37\ 25,19 & 37\ 30,77 & 36\ 39,79
\end{array}\]
wurden nach den Formeln:
\[ \operatorname{tang}\delta = \operatorname{tang} i \sin(\alpha-\AscendingNode) \]
\[ \operatorname{tang}(v-\AscendingNode) = \frac{\operatorname{tang}(\alpha-\AscendingNode)}{\cos i} \]
die heliocentrischen Breiten wie folgt, gefunden:
\[\begin{array}{@{}r@{\;}r|r|r|r|r|r@{}}
\delta={}&+33^{0}\,41'6'',30 & +10^{0}\,27'5'',72 & -33^{0}\,40'33'',54 & +28^{0}\,14'54'',88 & +27^{0}\,20'33'',96 & -4^{0}\,53'0'',48
\end{array}\]
sowie die Längen in der Bahn:
\[\begin{array}{@{}r@{\;}r|r|r|r|r|r@{}}
v={}&275^{0}\,3'20'',76 & 333^{0}\,51'47'',21 & 69^{0}\,53'48'',50 & 228^{0}\,52'49'',17 & 298^{0}\,31'39'',70 & 1^{0}\,4'56'',77.
\end{array}\]
\par
Indem man des Weitern die Perihellänge und den Winkel \(\varphi\) nach den obigen Störungswerthen, wie folgt, ansetzt:
\[\begin{array}{@{}r@{\;}r|r|r|r|r|r@{}}
\tilde\omega={}&121^{0}\,7'53'',99 & 6'38'',57 & 3'17'',76 & 0'53'',65 & 120^{0}\,59'36'',29 & 45'59'',50
\\
\varphi={}&14\ 12\ 47,43 & 11\ 33,15 & 10\ 58,62 & 10\ 10,05 & 14\ \ 8\ 10,26 & 1\ 34,91
\end{array}\]
folgen die wahren und die mittlern Anomalien, sowie die mittlern Längen:
\[\begin{array}{@{}r@{\;}r|r|r|r|r|r@{}}
v-\tilde\omega={}&153^{0}\,55'26'',77 & 212^{0}\,45'8'',64 & 308^{0}\,50'30'',74 & 104^{0}\,51'55'',52 & 177^{0}\,32'3'',41 & 240^{0}\,18'57'',27
\\
M={}&139\ 10\ 15,32 & 230\ 38\ 18,77 & 328\ 18\ 11,43 & 79\ 50\ 57,16 & 176\ \ 3\ 51,83 & 266\ 36\ \ 8,99
\\
L={}&260\ 18\ \ 9,31 & 351\ 44\ 57,34 & 89\ 21\ 29,19 & 200\ 51\ 50,81 & 297\ \ 3\ 28,12 & 27\ 22\ \ 8,49.
\end{array}\]
\par
Bringt man an die letztern Werthe die Störungsgrössen \(\delta\int n\,\mathrm{d}t+\delta\varepsilon\) mit umgekehrtem Zeichen an, so erhält man die mittlern Längen von den Störungen befreit, wie folgt:
\[\begin{array}{@{}r@{\;}r|r|r|r|r|r@{}}
L_0={}&260^{0}\,18'4'',58 & 351^{0}\,45'29'',31 & 89^{0}\,26'34'',76 & 201^{0}\,0'49'',51 & 297^{0}\,11'59'',74 & 27^{0}\,43'4'',39
\end{array}\]
und leitet man endlich aus diesen Werthen die Epoche der mittlern Länge für \(1803.156^{\mathrm d}\) ab, indem man die mittlere Bewegung, angenommen zu \(770'',7332\) multiplicirt mit der Zwischenzeit davon in Abzug bringt, so kommt
\[\begin{array}{@{}r@{\;}r|r|r|r|r|r@{}}
(\varepsilon_0)={}&254^{0}\,56'41'',51 & 56'38'',10 & 56'47'',58 & 56'42'',41 & 56'50'',65 & 56'47'',64.
\end{array}\]
\par
\SourceLeaf{487}Diese sechs Werte müssten einander gleich sein, wenn unsere Näherungswerthe für die Normalelemente \(\tilde\omega\), \(\varphi\) und \(n\) die richtigen wären; aus ihren gegenseitigen Abweichungen können nun die Correctionen für die Perihellänge, den Winkel \(\varphi\), die mittlere Bewegung und der wahre Werth für die Epoche der mittlern Länge abgeleitet werden. Sind \(\Delta\tilde\omega\) und \(\Delta\varphi\) die Correctionen für die Grössen \(\tilde\omega\) und \(\varphi\), so ist der corrigirte Werth für die mittlere Länge in jeder der Oppositionen gleich (vgl. Band VI Seite 7 oben)
\[ L_0 + (m+1)\Delta\tilde\omega + n\Delta\varphi, \]
wo die Coefficienten \(m\) und \(n\) dieselbe Bedeutung haben, wie a. a. O. Ist ferner \(\Delta n\) die Correction für die mittlere Bewegung, so wird die corrigirte Epoche der mittlern Länge nach Abzug der Störungen
\[ \varepsilon_0 = (\varepsilon_0) + (m+1)\Delta\tilde\omega + n\Delta\varphi - t\Delta n \tag*{A)}\]
wo \(t\) die seit \(1803.156^{\mathrm d}\) verflossene Zeit und \((\varepsilon_0)\) die oben aus den Beobachtungen gerechneten Werte bedeutet.
\par
Aus diesen sechs Gleichungen kann man zunächst die Epoche eliminiren, indem man das Mittel aus ihnen nimmt, welches den corrigirten Werth der Epoche in der Form:
\[ \varepsilon_0 = \tfrac{1}{6}\Sigma(\varepsilon_0) + \tfrac{1}{6}\Sigma(m+1).\Delta\tilde\omega + \tfrac{1}{6}\Sigma n.\Delta\varphi - \tfrac{1}{6}\Sigma t.\Delta n \tag*{B)}\]
ergibt. Indem man dann jede der vorstehenden Gleichungen A) von diesem Mittel B) abzieht, erhält man sechs Gleichungen, welche nur die Unbekannten \(\Delta\tilde\omega\), \(\Delta\varphi\), \(\Delta n\) enthalten und welche man nach der Methode der kleinsten Quadrate auflöst. Im vorliegenden Falle ergibt die numerische Rechnung, dass \(\Delta\tilde\omega\) und \(\Delta\varphi\) äusserst klein sind, während
\[ \Delta n = +0'',0040 \]
folgt, womit die mittlere Bewegung wird:
\[ n = 770'',7372. \]
\par
Für die Epoche ergibt dann Gleichung B) den Werth:
\[ \varepsilon_0 = 254^{0}\,56'40'',14. \]
\par
Corrigirt man nun die oben gefundenen sechs Werthe von \((\varepsilon_0)\) durch Anbringung der eben gefundenen Correctionen nach den Gleichungen A), so werden diese jetzt:
\[\begin{array}{@{}r@{\;}r|r|r|r|r|r@{}}
\varepsilon_0={}&254^{0}\,56'41'',41 & 56'36'',35 & 56'44'',07 & 56'36'',88 & 56'43'',38 & 56'38'',74
\end{array}\]
so dass also nur die folgenden sehr geringfügigen Abweichungen der mittlern Längen übrig bleiben (vgl. Seite 418 dieses Bandes und Band VI Seite 320)
\[\begin{array}{@{}r@{\;}r|r|r|r|r|r@{}}
&+1'',27 & -3'',79 & +3'',93 & -3'',26 & +3'',24 & -1'',40.
\end{array}\]
\par
Zur Verbesserung der Werte von \(i\) und \(\AscendingNode\) verfährt man, wie folgt: aus den eben gefundenen Werthen von \(n\) ergibt sich für die verschiedenen Oppositionen:
\[\begin{array}{@{}r@{\;}r|r|r|r|r|r@{}}
n={}&770'',7255 & 770'',3928 & 770'',1373 & 770'',5010 & 770'',1580 & 768''\IIIfourLowComma2969
\end{array}\]
hieraus
\[\begin{array}{@{}r@{\;}r|r|r|r|r|r@{}}
p={}&0{,}415\,0673 & 0{,}415\,2715 & 0{,}415\,4043 & 0{,}415\,3194 & 0{,}415\,5755 & 0{,}416\,6936
\end{array}\]
und
\[\begin{array}{@{}r@{\;}r|r|r|r|r|r@{}}
r={}&0{,}523\,2726 & 0{,}515\,5612 & 0{,}353\,3230 & 0{,}449\,2234 & 0{,}537\,0545 & 0{,}472\,2233
\end{array}\]
\SourceLeaf{488}
weiter nach der Formel \(\sin(\beta-\gamma) = \dfrac{R\sin\beta}{r}\):
\[\begin{array}{@{}r@{\;}r|r|r|r|r|r@{}}
\beta-\gamma={}&12^{0}\,45'32'',02 & 4^{0}\,34'41'',98 & 20^{0}\,50'28'',02 & 13^{0}\,56'26'',89 & 10^{0}\,23'37'',27 & 2^{0}\,29'4'',78
\\
\gamma={}&+33\ 41\ \ 3,98 & +10\ 27\ \ 7,82 & -33\ 40\ 26,88 & +28\ 14\ 58,71 & +27\ 20\ 16,43 & -4\ 53\ \ 5,32
\end{array}\]
und hiemit endlich die Differenz
\[\begin{array}{@{}r@{\;}r|r|r|r|r|r@{}}
\gamma-\delta={}&-2'',32 & +2'',10 & +6'',66 & +3'',83 & -17'',53 & -4'',84.
\end{array}\]
\par
Nach Art. 6 der Disquisitio erhält man dann für jede Opposition eine Gleichung der Form:
\[ \IIIfourLowC\Delta\AscendingNode + f\Delta i = \gamma-\delta \]
wo die Grössen \(a\) und \(b\) fortgelassen sind und die \(c\) und \(f\) nahe dieselben numerischen Werthe haben, wie a. a. O.; unsere Gleichungen liefern uns die Correctionen, wenn wir die V. Opposition bei Seite lassen, da sie schlecht beobachtet war:
\[ \Delta\AscendingNode = -2'',5 \quad\text{und}\quad \Delta i = -1'', \]
womit
\[ \AscendingNode = 172^{0}\,28'29'',5 \quad\text{und}\quad i = 34^{0}\,37'46'' \]
wird. Nach Anbringung der hieraus folgenden Correctionen an die obigen Werte von \(\gamma\) bleiben folgende Unterschiede \(\gamma-\delta\) in den geocentrischen Breiten übrig (vgl. Seite 418 dieses Bandes und Band VI Seite 320):
\[\begin{array}{@{}r@{\;}r|r|r|r|r|r@{}}
&-1'',00 & +4'',11 & +5'',99 & +3'',98 & -15'',86 & -3'',34.
\end{array}\]
\par
Dass gefundene System der Normalelemente ist also das folgende:
\[\begin{aligned}
n&=770'',7372\\
\tilde\omega&=121^{0}\,7'54''\\
\varphi&=14\ 12\ 51\\
\log a&=0{,}442\,0669\\
\varepsilon&=254^{0}\,56'40'',14\\
i&=34\ 37\ 46\\
\AscendingNode&=172\ 28\ 29,5.]
\end{aligned}\]
\par\medskip\begin{center}\rule{.45\linewidth}{.3pt}\end{center}
\par\endgroup
\end{landscape}
